% Einheit 3 von 5 – Brüche multiplizieren und dividieren (bank/bruchrechnung/e3.jsonl, zusammenbau v0.1)
%% TODO Zweigzeile Teil 1: was hier gelernt wird (Satz in Schülersprache aus dem Einheitstitel) – Einheit 3
\einheitenkopf[e3]{Einheit 3 von 5 · Brüche multiplizieren und dividieren}
\zweigzeile{ab Kl. 6 (am Gymnasium ab Kl. 5) · keine P10-Aufgabe}
\setcounter{aufgabe}{17}

%% TODO Titel: Ich-kann-Satz fehlt in der Bank (Platzhalter: Kettenname) – Nr. 18
%% TODO Anweisung: ein Satz über der ersten Teilaufgabe fehlt in der Bank – Nr. 18
\begin{aufgabe}{Von heißt mal}
\begin{teile}
\teil Unterstreiche „von“ und schreibe die Malaufgabe auf: Zwei Fünftel von 35 Kindern fahren mit dem Rad. \leerfeld $\cdot$ \leerfeld
\end{teile}
\end{aufgabe}

%% TODO \verfahren{…}: Verfahrensüberschrift in Sachsprache fehlt in der Bank (2.3 b)
%% TODO Titel: Ich-kann-Satz fehlt in der Bank (Platzhalter: Kettenname) – Nr. 19
%% TODO Anweisung: ein Satz über der ersten Teilaufgabe fehlt in der Bank – Nr. 19
\begin{aufgabe}{Multiplizieren}
%% TODO \rechenplatz{n}: Schrittzahl des Grundfalls fehlt in der Bank (nur bei fester Schreibform, 2.3 b)
\begin{teile}
\teil „$\frac{1}{4}$ von 28 Kindern“ – heißt „von“ hier mal? \janein
\teil $\frac{2}{9} \cdot 4$
\teil $\frac{2}{5} \cdot \frac{3}{7}$
\teil $\frac{3}{8} \cdot \frac{4}{9}$
\teil $1\frac{1}{2} \cdot \frac{4}{5}$
\teil In einer Dose liegen 15 Murmeln: 7 grün, 5 weiß, 3 schwarz. Drei werden nacheinander ohne Zurücklegen gezogen. Tom behauptet: „Dreimal grün ist doppelt so wahrscheinlich wie dreimal weiß.“ Widerlege die Behauptung mit einer Rechnung. \hfill (P10 2019 OS)
\teil Sechs Karten liegen verdeckt, eine davon ist die Gewinnkarte. Anna zieht zuerst eine Karte, danach Ben eine der übrigen fünf. Anna zieht keine Gewinnkarte mit der Wahrscheinlichkeit $\frac{5}{6}$, danach zieht Ben sie mit $\frac{1}{5}$. Wie groß ist die Wahrscheinlichkeit, dass Ben die Gewinnkarte zieht? \hfill (P10 2024 OS)
\end{teile}
\end{aufgabe}

%% TODO \verfahren{…}: Verfahrensüberschrift in Sachsprache fehlt in der Bank (2.3 b)
%% TODO Titel: Ich-kann-Satz fehlt in der Bank (Platzhalter: Kettenname) – Nr. 20
%% TODO Anweisung: ein Satz über der ersten Teilaufgabe fehlt in der Bank – Nr. 20
\begin{aufgabe}{Dividieren}
%% TODO \rechenplatz{n}: Schrittzahl des Grundfalls fehlt in der Bank (nur bei fester Schreibform, 2.3 b)
\begin{teile}
\teil $\frac{2}{3} : 5$
\teil $6 : \frac{1}{4}$ – wie oft passt $\frac{1}{4}$ in 6?
\teil $\frac{1}{6} : \frac{3}{5}$
\teil $\frac{5}{8} : \frac{15}{16}$
\teil Ein Kanister enthält 9 Liter Apfelsaft. Er wird in Flaschen zu je $\frac{3}{4}$ Liter abgefüllt. Wie viele Flaschen werden voll?
\end{teile}
\end{aufgabe}

%% TODO Titel: Ich-kann-Satz fehlt in der Bank (Platzhalter: Kettenname) – Nr. 21
%% TODO Anweisung: ein Satz über der ersten Teilaufgabe fehlt in der Bank – Nr. 21
\begin{aufgabe}{Multiplizieren – Fehler finden, Anwendung, Darstellungswechsel}
\begin{teile}
\teil Emil rechnet: $\frac{2}{3} \cdot \frac{4}{5} = \frac{2}{3} \cdot \frac{5}{4} = \frac{10}{12}$. Finde den Fehler.
\teil Ein Rezept für 4 Personen braucht $\frac{3}{4}$ Liter Milch. Du kochst für 6 Personen. Wie viel Milch brauchst du?
\teil Der Streifen hat 12 Teile. Färbe $\frac{1}{2}$ von $\frac{2}{3}$ des Streifens und gib den Anteil an.

\bruchrechteck{0}{12}
\leerfeld
\end{teile}
\end{aufgabe}

%% TODO Titel: Ich-kann-Satz fehlt in der Bank (Platzhalter: Kettenname) – Nr. 22
%% TODO Anweisung: ein Satz über der ersten Teilaufgabe fehlt in der Bank – Nr. 22
\begin{aufgabe}{Dividieren – Fehler finden, Begründen}
\begin{teile}
\teil Lars rechnet: $\frac{5}{8} : 4 = \frac{5}{2}$. Finde den Fehler.
\teil Warum ist $9 : \frac{1}{2}$ dasselbe wie $9 \cdot 2$?
\end{teile}
\end{aufgabe}
